\chapter{Аудит
\texorpdfstring{$\operatorname{Spin}^h$}{Spin-h}-связывания ориентации и
семейного переноса}

Предыдущий гейт доказал, что в текущем $H_{15}/M_{35}$ отсутствует
полуцелое семейное представление. Единственным оставшимся условием
переоткрытия была попытка использовать уже существующий ориентационный
дублет локального переноса и связать его с семейным $SO(3)$ языком
$\operatorname{Spin}^h$. Настоящая глава проверяет эту возможность до
введения нового действия.

Гейт относится к аудиту возможной новой версии. Он не отменяет
представительный запрет версии V и не объявляет дополнительный сектор
разрешённым.

\section{Точный смысл структуры
\texorpdfstring{$\operatorname{Spin}^h$}{Spin-h}}

Группа определяется как
\begin{equation}
 \operatorname{Spin}^h(n)=
 \frac{\operatorname{Spin}(n)\times Sp(1)}
 {\langle(-1,-1)\rangle},
 \qquad Sp(1)\simeq SU(2).
 \label{eq:v5-spinh-definition}
\end{equation}
Она является кватернионным аналогом $\operatorname{Spin}^c$. Представление
факторгруппы допустимо только тогда, когда диагональная пара
$(-1,-1)$ действует тождественно.

Геометрически структура несёт вспомогательное главное $SO(3)$-расслоение
$Q$ и требует согласования
\begin{equation}
 w_2(Q)=w_2(TM).
 \label{eq:v5-spinh-w2-condition}
\end{equation}
То есть $\operatorname{Spin}^h$ не превращает произвольный семейный
$SO(3)$ в спинор автоматически. Он требует отдельного полуцелого действия
$Sp(1)$, центральный знак которого компенсирует знак пространственного
спинора.

Это согласуется с математическим определением
$\operatorname{Spin}^r(n)=(\operatorname{Spin}(n)\times
\operatorname{Spin}(r))/\mathbb Z_2$ и с физическим чтением
$\operatorname{Spin}^h$-спинора как спинора с полуцелым изоспином,
связанного с $SO(3)$-полем.

\section{Можно ли поместить два действия на одну
\texorpdfstring{$\mathbb C^2$}{C2}}

Ориентационный множитель локального переноса имеет размерность два:
\begin{equation}
 \mathcal O=\mathbb C^2_{\rm orient}.
\end{equation}
Предположим наиболее благоприятно, что на нём уже действует
фундаментальное пространственно-спинорное $SU(2)$. Чтобы тот же носитель
реализовал $\operatorname{Spin}^h$, на нём должно существовать второе
независимое фундаментальное семейное действие, коммутирующее с первым,
поскольку до факторизации имеется прямая группа
$\operatorname{Spin}(n)\times SU(2)_F$.

Но фундаментальное представление $SU(2)$ на $\mathbb C^2$ неприводимо.
Его коммутант равен
\begin{equation}
 \operatorname{End}_{SU(2)}(\mathbb C^2)=\mathbb C I_2.
 \label{eq:v5-spinh-c2-commutant}
\end{equation}
В скалярном коммутанте нет трёх линейно независимых бесследовых
генераторов второго $\mathfrak{su}(2)$. Машинный аудит строит линейную
систему для общего элемента в базисе
$\{I,\sigma_x,\sigma_y,\sigma_z\}$ и получает одномерный коммутант.

\begin{theorem}[Запрет двойной роли одной ориентационной пары]
Одна комплексная двойка не может одновременно нести два независимых
нетривиальных коммутирующих действия $SU(2)$: пространственно-спинорное и
семейное. Следовательно, существующий ориентационный дублет не является
скрытым $\operatorname{Spin}^h$-модулем.
\end{theorem}

\section{Минимальный настоящий
\texorpdfstring{$\operatorname{Spin}^h$}{Spin-h}-модуль}

Минимальное прямое решение имеет вид
\begin{equation}
 \mathcal S_h=
 \mathbb C^2_{\rm spin}\otimes\mathbb C^2_F.
 \label{eq:v5-spinh-minimal-module}
\end{equation}
Первый $SU(2)$ действует как $U\otimes I_2$, второй --- как
$I_2\otimes V$. Эти действия коммутируют, а диагональный центральный
элемент даёт
\begin{equation}
 (-I_2)\otimes(-I_2)=+I_4.
 \label{eq:v5-spinh-center-cancellation}
\end{equation}
Поэтому четырёхмерное комплексное представление действительно спускается
на факторгруппу \eqref{eq:v5-spinh-definition}.

Однако в текущем проекте имеется только одна внешняя ориентационная
$\mathbb C^2$. Вторая, семейная $\mathbb C^2_F$, не содержится ни в
$H_{15}$, ни в $M_{35}$. Для локального переноса размерность меняется как
\begin{equation}
 2\cdot300=600
 \quad\longrightarrow\quad
 2\cdot2\cdot300=1200.
 \label{eq:v5-spinh-dimension-budget}
\end{equation}
Это не просто другое обозначение тех же состояний, а новый внутренний
дублет.

\section{Почему диагональное отождествление не спасает размерность}

Можно попытаться объявить пространственный и семейный $SU(2)$ одной
диагональной группой. Тогда исчезает требование двух коммутирующих
действий, но исчезает и исходная архитектура
$\operatorname{Spin}(n)\times SU(2)_F$: пространственные вращения начинают
одновременно вращать семейную ось.

Кроме того, фундаментальная $\mathbb C^2$ диагональной группы сама имеет
центральное действие $-I_2$, тогда как элемент $(-1,-1)$ должен стать
единицей факторгруппы. Простое переименование одной матрицы Паули не
реализует представление $\operatorname{Spin}^h$ и не воспроизводит
компенсацию \eqref{eq:v5-spinh-center-cancellation}.

Физически такое отождествление связало бы обычные пространственные
вращения с поколенческими преобразованиями. Ни Том IV, ни родитель Мориты
не содержат подобной спин-семейной блокировки.

\section{Топологическая проверка}

Текущий носитель $\mathbb{RP}^3\times S^1$ уже допускает обычные
спин-структуры. Поэтому $\operatorname{Spin}^h$ здесь не требуется для
устранения препятствия касательного расслоения. Он мог бы быть новой
совместной структурой, но тогда нужно отдельно построить вспомогательное
$SO(3)$-расслоение $Q$ и доказать
\eqref{eq:v5-spinh-w2-condition}.

Проекторное поле $P=nn^T$ задаёт отображение в $\mathbb{RP}^2$, но само по
себе ещё не задаёт главное расслоение $Q$, его переходные функции и
$w_2(Q)$. Поэтому топологическое условие $\operatorname{Spin}^h$ не
следует из уже вычисленного заряда ежа.

\section{След и аномалии}

Новый множитель $\mathbb C^2_F$ не нормируется существующим следом
$M_{35}$. Потребуется расширенный родительский след и повторный вывод всех
относительных весов. Если полуцелый семейный дублет назначить всем каналам
$H_{15}$, снова возникают пятнадцать хиральных дублетов. Простое название
$\operatorname{Spin}^h$ не отменяет глобальный аномальный тест; требуется
полный бордизмовый и фермионно-детерминантный аудит новой модели.

\begin{nogocriterion}[Запрет скрытого $\operatorname{Spin}^h$-спасения]
Ориентационная $\mathbb C^2$ не может одновременно играть роль
пространственного и независимого семейного дублета. Настоящий
$\operatorname{Spin}^h$-модуль требует дополнительного множителя
$\mathbb C^2_F$, нового следа и отдельно выведенного $SO(3)$-расслоения.
Следовательно, он не переоткрывает версию V без расширения модели.
\end{nogocriterion}

\section{Вердикт}

$\operatorname{Spin}^h$ оказался не пустой метафорой, а точным языком того,
какой могла бы быть новая архитектура. Но он одновременно показал, что
искомого механизма нет в текущем наборе состояний: для компенсации двух
центральных знаков нужны две независимые двойки, а имеется одна.

\begin{protocol}
\begin{enumerate}
 \item математический язык $\operatorname{Spin}^h$:
 \textbf{применим};
 \item одна ориентационная $\mathbb C^2$ как два независимых дублета:
 \textbf{невозможна};
 \item минимальный модуль: \textbf{$\mathbb C^2\otimes\mathbb C^2$};
 \item дополнительный семейный дублет в $H_{15}/M_{35}$:
 \textbf{отсутствует};
 \item условие $w_2(Q)=w_2(TM)$: \textbf{не выведено};
 \item нормировка новым следом: \textbf{не выведена};
 \item глобальная аномалия: \textbf{требует нового аудита};
 \item переоткрытие без новых модулей: \textbf{отклонено};
 \item маршрут версии V: \textbf{остаётся закрытым};
 \item автоматический следующий гейт: \textbf{отсутствует}.
\end{enumerate}
\end{protocol}

Машинный протокол сохранён в
\path{s2t_v5_spinh_orientation_family_locking_reopening_gate_results.json}.